\section{Introducción}
\label{sec:introduccion}
% Planteo del problema, justificación, objetivos (general y específicos), estructura del informe.
% Tiene que ser más larga que el resumen.

La suma de números reales es conmutativa y asociativa. Esas dos propiedades permiten escribir una
sumatoria como $\sum_{k=1}^{N} x_k$ sin aclarar en qué orden se recorren los términos ni cómo se
agrupan los paréntesis: cualquier camino conduce al mismo valor. La aritmética que ejecuta una
computadora no ofrece esa garantía. Cada resultado intermedio debe caber en una cantidad fija de
bits, así que casi todas las operaciones redondean. Ese redondeo depende de los valores que llegan
a cada paso. Dos programas que calculan la misma sumatoria con distinto orden o distinto
agrupamiento son entonces dos algoritmos diferentes, y pueden devolver números diferentes
(Goldberg, 1991).

Esa distancia entre la propiedad algebraica y su implementación aparece en situaciones corrientes de
programación. Un mismo código puede dar resultados distintos al paralelizarse, al cambiar de
biblioteca o al reordenar un bucle. Una suma de muchos términos pequeños puede perder cifras
significativas sin que nada en el programa lo señale. El problema no se anuncia con un mensaje de
error: el cálculo termina, devuelve un número con quince dígitos y algunos de esos dígitos son
incorrectos. Detectarlo exige una referencia contra la cual medir el error y un modelo de los
mecanismos que lo producen. Ambos son objeto del análisis numérico del redondeo (Higham, 2002).

Este trabajo estudia ese fenómeno sobre tres sucesiones elegidas para hacerlo medible. La primera,
$a_N = \sum_{k=1}^{N} (1 + b^{k} - b^{k})$, vale exactamente $N$ en $\mathbb{R}$ para cualquier base
$b$, de modo que toda desviación observada proviene del cálculo y no del planteo. La segunda,
$b_N = \sum_{k=1}^{N} 1/(k(k+1))$, tiene forma cerrada $N/(N+1)$ y admite varias formas de
recorrerse y de escribirse, lo que permite comparar algoritmos contra un valor de referencia. La
tercera, $c_N = \sum_{k=1}^{N} 1/(\sqrt{k^{2}+1} + k)$, no tiene forma cerrada conocida pero sí una
representación algebraicamente equivalente, $\sqrt{k^{2}+1} - k$, expuesta a la cancelación entre
cifras próximas. Las tres comparten una condición: el resultado correcto se conoce de antemano o se
puede acotar, y sin esa condición el error no sería observable.

El objetivo general es determinar en qué condiciones las propiedades algebraicas de la suma dejan de
valer en la aritmética de una computadora, cuánto se aparta el resultado del valor exacto en cada
caso y qué explica ese apartamiento. De ahí se desprenden cinco objetivos específicos:

\begin{enumerate}
  \item Identificar a partir de qué tamaño de los sumandos la agrupación de los paréntesis altera el
  resultado de una suma, y contrastar el comportamiento del entero de precisión arbitraria de Python
  con el del punto flotante de doble precisión y con el del entero de ancho fijo.
  \item Medir cómo depende el error de una suma acumulada del orden en que se recorren los términos,
  comparando cuatro estrategias de acumulación sobre la misma serie, entre ellas la suma compensada
  de Kahan.
  \item Cuantificar la discrepancia entre representaciones algebraicamente equivalentes de una misma
  sumatoria y establecer cuándo la cancelación entre cantidades próximas degrada el resultado y
  cuándo resulta inofensiva.
  \item Contrastar cada medición con las cotas de error que predice el análisis teórico, y evaluar
  si describen el comportamiento observado o lo sobreestiman.
  \item Producir un conjunto de programas que genere todas las figuras y tablas del informe, de modo
  que cada dato presentado sea reproducible a partir del código.
\end{enumerate}

El informe se organiza en cinco secciones además de esta, más la bibliografía. La sección \ref{sec:marco} reúne los
conceptos necesarios para leer los resultados: la representación en punto flotante y sus dos
constantes de resolución, los enteros de precisión arbitraria, el estado de las propiedades
conmutativa y asociativa en la aritmética de máquina, la medida de error empleada, la suma de Kahan,
las sumas telescópicas y la cancelación catastrófica. Allí quedan enunciadas las cotas que después
se comparan contra los datos. La sección \ref{sec:metodologia} describe el entorno, el diseño de
cada experimento, los criterios de análisis, el tratamiento de la aleatoriedad y la ubicación del
código. La sección
\ref{sec:resultados} presenta las mediciones agrupadas por fenómeno y no por experimento: la pérdida
de la asociatividad, la dependencia del orden de suma y el comportamiento de representaciones
equivalentes. La sección \ref{sec:discusion} interpreta esas mediciones, las compara con las cotas
teóricas, examina dos resultados que contradicen la expectativa inicial y enuncia las limitaciones
del estudio. La sección \ref{sec:conclusiones} sintetiza las respuestas a los objetivos y señala
líneas de trabajo posteriores.
